\scanpage{110}

\[
T\longrightarrow Y
\]
\uline{and} of a lift to $T\longrightarrow X$, or what amounts to the same, of a section of $X_T=X\times_YT$ over $T$. This being so,

\noindent
\begin{minipage}[t]{0.48\linewidth}
\vspace*{\baselineskip}
\diagram{\begin{tikzcd}
X\arrow[d] & X_T\arrow[l]\arrow[d] & \\
Y & T\arrow[l] & T\arrow[ul, "{\unsure{}}"{description,pos=0.18}]\arrow[l,phantom,"=" description]
\end{tikzcd}}
% The question near the diagonal's lower-right end marks the possible second
% arrowhead there. The upper-left head is certain; the short lower-right return
% may be a tail return or a second head. See the whole-diagram audit and doubt.
\end{minipage}\hfill
\begin{minipage}[t]{0.48\linewidth}
we shall have a canonical isomorphism (under the hypothesis of ``homotopy exact sequence'' \deleted{\unsure{)}} made)
\end{minipage}

\[
(12)\qquad \mathcal E_{X/Y}^T\simeq\Pi_1(X_T;T)\simeq\mathcal E_{X_T}^T
\]
and we shall have
\[
(13)\qquad 1\longrightarrow\mathcal E_{X/Y}^T\longrightarrow\mathcal E_X^T\longrightarrow\mathcal E_Y^T\longrightarrow1
\]
(This hypothesis will always be satisfied if $Y=\operatorname{Spec}K$, $K$ a field, if $X$ is geometrically $0$-connected over $K$)

More generally, if we have a universal homotopy exact sequence, but not necessarily with $\mathcal E_X^T\longrightarrow\mathcal E_Y^T$ surjective,

\margin{109}
